\documentclass[11pt]{article}
\usepackage[T1]{fontenc}
\usepackage[utf8]{inputenc}
\usepackage{lmodern,microtype}
\usepackage[margin=1.08in]{geometry}
\usepackage{amsmath,amssymb,amsthm,mathtools}
\usepackage{xcolor,hyperref}
\hypersetup{colorlinks=true,linkcolor=blue!55!black,citecolor=blue!55!black,urlcolor=blue!60!black}
\setlength{\parindent}{0pt}
\setlength{\parskip}{0.48em}
\setlength{\emergencystretch}{2em}
\newtheorem{theorem}{Theorem}
\newtheorem{lemma}[theorem]{Lemma}
\theoremstyle{remark}
\newtheorem{remark}[theorem]{Remark}

\title{Umbilics of Analytic $W$-Congruences with
\texorpdfstring{$m\notin\mathbb N$}{m not in N}
Have \texorpdfstring{$A_\infty$}{A-infinity} Discriminant}
\author{[Author Name]}
\date{}

\begin{document}
\maketitle

\begin{abstract}
\textbf{Relation to the problem list.}
This manuscript addresses Question~10 (L10), ``Generic noninteger
$W$-singularity.''  It gives a full proof in the real-analytic,
nondegenerate, fixed-$m$ setting; in fact it proves the stronger statement
that every admissible germ in that stratum has type $A_\infty$, rather than
only a generic set of germs.

We study nondegenerate umbilics of real-analytic $W$-congruences in the
normalization of Craizer and Garcia.  For every admissible value
$m\notin\mathbb N$
of the normalized parameter $m$, we prove that the discriminant germ is
analytically right-equivalent to the square of one coordinate.  Thus the
umbilic has type $A_\infty$.  The proof converts the $W$-equation into an
eigenfunction equation $X\Delta=2\Delta$ for an analytic planar vector
field.  A nonresonant analytic invariant curve and two one-dimensional
order comparisons force $\Delta$ to have a double transverse factor.
\end{abstract}

\section{Introduction and statement}

In an affine chart, a line congruence may be represented by
\[
 f(x,y)=(x,y,0),
 \qquad
 \xi(x,y)=(P(x,y),Q(x,y),1),
\]
where $P$ and $Q$ are real analytic.  The extended Weingarten equation is
\begin{align}
 W(P,Q)={}&(Q_y-P_x)(Q_{xx}P_{yy}-P_{xx}Q_{yy})\notag\\
 &-2P_y(P_{xx}Q_{xy}-Q_{xx}P_{xy})
 +2Q_x(P_{xy}Q_{yy}-P_{yy}Q_{xy})=0.\label{eq:W}
\end{align}
Its discriminant is
\begin{equation}\label{eq:disc}
 \delta=(P_x-Q_y)^2+4P_yQ_x.
\end{equation}
At an umbilic at the origin one has
$P_x=Q_y$ and $P_y=Q_x=0$.  Subtracting the common linear part and making
the standard linear normalization, write
\[
 p_{ij}=\partial_x^i\partial_y^jP(0),
 \qquad
 q_{ij}=\partial_x^i\partial_y^jQ(0),
\]
and assume
\begin{equation}\label{eq:normalization}
 p_{20}=p_{02}=q_{11}=0,
 \qquad
 p_{11}=a\ne0,
 \qquad
 q_{02}\ne0,
 \qquad
 p_{11}\ne q_{02}.
\end{equation}
The normalized parameter is
\begin{equation}\label{eq:m}
 m=-\frac{q_{02}}{a}.
\end{equation}
The hypotheses exclude $m=0$ and $m=-1$.  A germ is said to have type
$A_\infty$ when its discriminant is analytically right-equivalent to
$v^2$.

\begin{theorem}\label{thm:main}
Let $(P,Q)$ be a real-analytic solution of \eqref{eq:W} with a
nondegenerate umbilic satisfying \eqref{eq:normalization}.  If the parameter
$m$ in \eqref{eq:m} is not a positive integer, then there are analytic
coordinates $(u,v)$ centered at the umbilic in which
\[
 \delta(u,v)=v^2.
\]
In particular, every admissible fixed-$m$ germ with $m\notin\mathbb N$ has type
$A_\infty$.
\end{theorem}

\section{Hodograph form of the \texorpdfstring{$W$}{W}-equation}

Because $m\ne1$, the first differentiated $W$-equation gives $q_{20}=0$.
Set
\begin{equation}\label{eq:hodograph}
 A=P_y,
 \qquad
 B=\frac{P_x-Q_y}{2},
 \qquad
 -Q_x=F(A,B),
 \qquad
 \frac{P_x+Q_y}{2}=G(A,B).
\end{equation}
At the origin,
\[
 A=ax+O(2),
 \qquad
 B=\frac{m+1}{2}ay+O(2),
\]
so $(A,B)$ are analytic local coordinates.  The normalized two-jet gives
\[
 F=O(2),
 \qquad
 G=\frac{1-m}{1+m}B+O(2).
\]
Define
\begin{align}
 \Delta&=B^2-AF,\label{eq:Delta}\\
 X&=A(1+G_B)\partial_A+(B-AG_A)\partial_B,\label{eq:X}\\
 E&=(1-G_B)F-(1+G_B)AF_A
 +(AG_A-B)F_B-2BG_A.\label{eq:E}
\end{align}

\begin{lemma}\label{lem:identities}
Let $J=\det(\partial(A,B)/\partial(x,y))$.  Then
\begin{equation}\label{eq:identities}
 W=2JE,
 \qquad
 (X-2)\Delta=AE,
 \qquad
 \delta=4\Delta.
\end{equation}
Consequently, every analytic $W$-germ satisfies
\begin{equation}\label{eq:eigenfunction}
 X\Delta=2\Delta.
\end{equation}
Moreover,
\begin{equation}\label{eq:linear}
 DX(0)=\operatorname{diag}(\alpha,1),
 \qquad
 \alpha=\frac{2}{m+1}\ne0,
 \qquad
 \Delta=B^2+O(3).
\end{equation}
\end{lemma}

\begin{proof}
The derivative relations in \eqref{eq:hodograph} are
\[
 P_x=B+G,
 \quad P_y=A,
 \quad Q_x=-F,
 \quad Q_y=G-B.
\]
Substitution of these four identities and their first derivatives into
\eqref{eq:W} gives $W=2JE$.  Applying the vector field \eqref{eq:X} to
$B^2-AF$ and collecting terms gives $(X-2)\Delta=AE$.  The last identity
in \eqref{eq:identities} follows immediately from \eqref{eq:disc}.
Since $J(0)=\frac{m+1}{2}a^2\ne0$, the equation $W=0$ implies $E=0$;
the remaining assertions follow from the normalized linear terms.
\end{proof}

\section{An analytic invariant curve}

\begin{lemma}\label{lem:curve}
Under the hypotheses of Theorem~\ref{thm:main}, the vector field $X$ has
an analytic invariant curve through the origin tangent to the $A$-axis.
\end{lemma}

\begin{proof}
Write
\[
 X=(\alpha A+R(A,B))\partial_A+(B+S(A,B))\partial_B,
 \qquad R,S=O(2).
\]
Seek the invariant curve as $B=\phi(A)$ with $\phi=O(A^2)$.  Its
invariance equation is
\begin{equation}\label{eq:invariance}
 \phi'(A)\bigl(\alpha A+R(A,\phi(A))\bigr)
 =\phi(A)+S(A,\phi(A)).
\end{equation}
Writing $\phi(A)=\sum_{n\ge2}\phi_nA^n$ gives the triangular recursion
\begin{equation}\label{eq:recursion}
 (n\alpha-1)\phi_n
 =\Psi_n(\phi_2,\ldots,\phi_{n-1};R,S).
\end{equation}
The only possible zero denominator satisfies
\[
 n\alpha=1
 \quad\Longleftrightarrow\quad
 m=2n-1,
\]
which is a positive integer.  It is therefore excluded.  For fixed
$\alpha\ne0$, the nonzero denominators in \eqref{eq:recursion} grow
linearly in $n$.  The standard majorant argument for analytic
Briot--Bouquet equations makes the recursively defined series convergent;
see \cite[Chapter~2]{IlyashenkoYakovenko}.  Hence the invariant curve is
analytic.
\end{proof}

Choose analytic coordinates $(s,t)$ for which this curve is $\{t=0\}$,
with $s$ tangent to the $A$-axis and $t$ transverse to it.  Invariance gives
\begin{equation}\label{eq:Xst}
 X=v(s,t)\partial_s+t w(s,t)\partial_t,
 \qquad
 v(s,0)=\alpha s+O(s^2),
 \qquad
 w(0,0)=1.
\end{equation}

\section{Double divisibility of the discriminant}

\begin{proof}[Proof of Theorem~\ref{thm:main}]
Let $f(s)=\Delta(s,0)$.  Restricting \eqref{eq:eigenfunction} to the
invariant curve gives
\begin{equation}\label{eq:firstrestriction}
 v(s,0)f'(s)=2f(s).
\end{equation}
If $f$ is not identically zero and has leading term $cs^n$, then
$n\ge3$ by \eqref{eq:linear}.  Comparing leading terms in
\eqref{eq:firstrestriction} yields
\[
 \alpha n=2
 \quad\Longleftrightarrow\quad
 m=n-1,
\]
contrary to the hypothesis.  Thus $f\equiv0$, and analytic division gives
$\Delta=tH$.

Because $Xt=tw$, the equation $X(tH)=2tH$ becomes
\[
 XH+(w-2)H=0.
\]
Let $h(s)=H(s,0)$.  Restriction to the invariant curve yields
\begin{equation}\label{eq:secondrestriction}
 v(s,0)h'(s)+(w(s,0)-2)h(s)=0.
\end{equation}
If $h\not\equiv0$ has leading term $ds^n$, where $n\ge0$, comparison of
leading terms gives $\alpha n+1=2$.  The case $n=0$ is impossible, while
for $n\ge1$ this is equivalent to
\[
 m=2n-1,
\]
again a positive integer.  Therefore $h\equiv0$ and a second analytic
division gives
\begin{equation}\label{eq:factor}
 \Delta=t^2U(s,t).
\end{equation}
The quadratic term in \eqref{eq:linear} and transversality of $t$ imply
$U(0,0)>0$.  Hence
\[
 \widetilde t=t\sqrt{U(s,t)}
\]
is an analytic coordinate.  By \eqref{eq:identities} and
\eqref{eq:factor},
\[
 \delta=4\widetilde t^{,2}.
\]
Rescaling $\widetilde t$ completes the analytic right equivalence to the
normal form $v^2$.
\end{proof}

\begin{remark}
The conclusion is relative to the fixed-$m$ analytic nondegenerate stratum.
It is stronger than a generic statement: every germ in that stratum has
type $A_\infty$.  The proof uses analyticity essentially in the invariant
curve, analytic division, and the square-root coordinate change.
\end{remark}

\begin{thebibliography}{9}

\bibitem{CraizerGarcia}
M.~Craizer and R.~A.~Garcia,
\emph{Singularities of Weingarten line congruences},
\href{https://arxiv.org/abs/2504.07289}{arXiv:2504.07289}.

\bibitem{IlyashenkoYakovenko}
Y.~Ilyashenko and S.~Yakovenko,
\emph{Lectures on Analytic Differential Equations},
Graduate Studies in Mathematics, vol.~86,
American Mathematical Society, Providence, 2008.

\end{thebibliography}
\end{document}
